\documentclass[10pt]{article}
% Shared LaTeX style for MIT 18.06 exam revision notes.
% Compile topic files with Tectonic, XeLaTeX, or LuaLaTeX.

\usepackage[a4paper,margin=0.72in]{geometry}
\usepackage{fontspec}
\usepackage{amsmath,amssymb,mathtools}
\usepackage{booktabs,array,tabularx}
\usepackage{enumitem}
\usepackage{titlesec}
\usepackage{fancyhdr}
\usepackage{microtype}
\usepackage{xcolor}

\setmainfont{Times New Roman}
\setsansfont{Arial}

\setlength{\parindent}{0pt}
\setlength{\parskip}{3.6pt}
\setlength{\abovedisplayskip}{6pt}
\setlength{\belowdisplayskip}{6pt}
\setlength{\abovedisplayshortskip}{4pt}
\setlength{\belowdisplayshortskip}{4pt}
\renewcommand{\arraystretch}{1.08}
\setlength{\tabcolsep}{4.5pt}

\titleformat{\section}
  {\normalfont\large\bfseries}
  {\thesection}
  {0.55em}
  {}
\titlespacing*{\section}{0pt}{10pt}{3pt}

\titleformat{\subsection}
  {\normalfont\normalsize\bfseries}
  {\thesubsection}
  {0.5em}
  {}
\titlespacing*{\subsection}{0pt}{7pt}{2pt}

\setlist[itemize]{leftmargin=1.35em,itemsep=1pt,topsep=2pt,parsep=0pt}
\setlist[enumerate]{leftmargin=1.55em,itemsep=1pt,topsep=2pt,parsep=0pt}

\pagestyle{fancy}
\fancyhf{}
\fancyfoot[C]{\scriptsize\color{black!45}MIT 18.06 Linear Algebra Exam Notes --- \thepage}
\renewcommand{\headrulewidth}{0pt}
\renewcommand{\footrulewidth}{0pt}

\newcommand{\notesubtitle}[1]{%
  {\centering\itshape #1\par}
  \vspace{2pt}
}

\newcommand{\source}[1]{%
  {\centering\small #1\par}
  \vspace{6pt}
}

\newcommand{\labelnote}[2]{%
  \par\smallskip
  \noindent\textbf{#1.}\quad #2
  \par\smallskip
}

\newcommand{\tightlabel}[1]{%
  \par\smallskip
  \noindent\textbf{#1.}\quad
}

\newcolumntype{Y}{>{\raggedright\arraybackslash}X}
\newcolumntype{L}[1]{>{\raggedright\arraybackslash}p{#1}}

\newenvironment{notetable}
  {\small\begin{tabularx}{\linewidth}}
  {\end{tabularx}}


\begin{document}

\begin{center}
{\LARGE 5.3 Inverses, Cramer's Rule, Cross Products, and Volume\par}
\end{center}
\notesubtitle{Concise exam revision notes for adjugates, determinant geometry, and orientation.}
\source{Based on handwritten MIT 18.06 revision notes.}

\section{Core Idea}

\labelnote{Core idea}{Determinants connect algebra and geometry: they decide invertibility, give formulas for inverse entries, and measure signed area or volume.}

For an invertible square matrix,
\[
A^{-1}=\frac{1}{\det A}C^{T},
\]
where \(C\) is the cofactor matrix. Equivalently,
\[
(\det A)I=AC^{T}.
\]

\section{Cramer's Rule}

For
\[
Ax=b,
\]
each variable is a determinant ratio:
\[
x_i=\frac{\det A_i(b)}{\det A},
\]
where \(A_i(b)\) is \(A\) with column \(i\) replaced by \(b\).

\begin{center}
\small
\begin{tabularx}{\linewidth}{@{}L{0.26\linewidth}Y@{}}
\toprule
Use Cramer's rule when & Avoid it when \\
\midrule
The matrix is small or symbolic. & The matrix is large and elimination is easier. \\
You need a formula for one variable. & You need all variables numerically. \\
\bottomrule
\end{tabularx}
\end{center}

\labelnote{Common mistake}{Cramer's rule requires \(\det A\ne0\). If \(\det A=0\), the determinant ratios do not apply.}

\section{Area from Determinants}

For two vectors in the plane,
\[
u=
\begin{bmatrix}
a\\
b
\end{bmatrix},
\qquad
v=
\begin{bmatrix}
c\\
d
\end{bmatrix},
\]
the signed area of the parallelogram is
\[
\det
\begin{bmatrix}
a&c\\
b&d
\end{bmatrix}
=ad-bc.
\]
The absolute value gives the geometric area:
\[
\text{area}=|ad-bc|.
\]

\section{Cross Product}

In \(\mathbb{R}^3\), the cross product is perpendicular to both input vectors:
\[
u\times v\perp u,
\qquad
u\times v\perp v.
\]
Using unit vectors,
\[
u\times v=
\det
\begin{bmatrix}
i&j&k\\
u_1&u_2&u_3\\
v_1&v_2&v_3
\end{bmatrix}.
\]
Its length is the parallelogram area:
\[
\|u\times v\|=\|u\|\,\|v\|\sin\theta.
\]

\labelnote{Exam use}{Use cross product when a problem asks for a normal vector to a plane in \(\mathbb{R}^3\).}

\section{Orientation}

The determinant sign records orientation:
\[
\det A>0
\]
preserves orientation, while
\[
\det A<0
\]
reverses orientation.

The scalar triple product gives signed volume:
\[
\det
\begin{bmatrix}
u_1&v_1&w_1\\
u_2&v_2&w_2\\
u_3&v_3&w_3
\end{bmatrix}
=(u\times v)\cdot w.
\]
If this determinant is zero, the three vectors lie in the same plane.

\section{Common Mistakes}

\begin{center}
\small
\begin{tabularx}{\linewidth}{@{}L{0.40\linewidth}Y@{}}
\toprule
Mistake & Correct exam habit \\
\midrule
Using Cramer's rule when \(\det A=0\). & First check invertibility. \\
Forgetting the transpose in the inverse formula. & Use \(A^{-1}=C^{T}/\det A\). \\
Dropping the sign of area too early. & Use signed area for orientation; absolute value for size. \\
Writing \(u\times v=v\times u\). & Cross product changes sign: \(v\times u=-(u\times v)\). \\
Using cross product outside \(\mathbb{R}^3\). & The standard cross product is a 3D tool. \\
\bottomrule
\end{tabularx}
\end{center}

\section{Final Exam Checklist}

\tightlabel{Checklist}
\begin{enumerate}
  \item Check whether \(\det A\ne0\) before using inverse formulas.
  \item Build the cofactor matrix \(C\).
  \item Use \(A^{-1}=C^{T}/\det A\).
  \item For Cramer's rule, replace one column at a time by \(b\).
  \item Use determinant sign for orientation.
  \item Use absolute determinant for area or volume size.
  \item Use \(u\times v\) for a normal vector in \(\mathbb{R}^3\).
  \item Check \(u\cdot(u\times v)=0\) and \(v\cdot(u\times v)=0\).
\end{enumerate}

\end{document}
